\documentclass[10pt,twocolumn]{article}
\usepackage[scaled]{helvet}
\renewcommand\familydefault{\sfdefault}
\usepackage[T1]{fontenc}
\usepackage[margin=0.7in]{geometry}
\usepackage{titlesec}
\usepackage{enumitem}
\usepackage{parskip}
\setlength{\columnsep}{0.24in}
\setlist[itemize]{leftmargin=1.2em, topsep=0.2em, itemsep=0.18em}
\titleformat{\section}{\large\bfseries}{\thesection}{1em}{}

\begin{document}
\title{\vspace{-1.6cm}AWS SageMaker}
\author{AWS ML Study Notes}
\date{}
\maketitle
\small

\section*{Overview}
SageMaker is AWS's managed machine learning platform for building, training, tuning, deploying, and governing ML models. It reduces the amount of infrastructure engineering needed to operationalize models.

Rather than giving you only raw compute, SageMaker gives purpose built ML abstractions such as training jobs, endpoints, pipelines, model registry integration, experiment tooling, and managed notebooks.

\section*{Core Concepts}
\begin{itemize}
\item SageMaker separates major concerns into specific services such as Studio, training, processing, batch transform, endpoints, feature related tooling, and pipeline orchestration.
\item A typical SageMaker workflow uses S3 for data and artifacts, IAM for execution roles, CloudWatch for logs and metrics, and VPC integration when private networking is required.
\item The value of SageMaker is not only managed infrastructure. It also standardizes packaging, deployment, lineage, and repeatability across teams.
\end{itemize}

\section*{How It Works}
\begin{itemize}
\item You prepare data, choose or build a container image, define the training configuration, and launch a managed job that provisions compute only for the duration needed.
\item The trained model artifact can then be registered, deployed to a real time endpoint, served asynchronously, used for batch transform, or promoted through an MLOps pipeline.
\item Operational events such as logs, metrics, model approval, and endpoint scaling are integrated into the wider AWS control plane.
\end{itemize}

\section*{AI and ML Relevance}
\begin{itemize}
\item SageMaker is useful when you want managed distributed training, hyperparameter tuning, controlled deployment patterns, or an internal ML platform without building every layer from scratch.
\item It is especially strong for teams that need repeatable ML workflows with governance, auditability, and tighter integration with AWS-native security and networking patterns.
\end{itemize}

\section*{Operational Notes}
\begin{itemize}
\item SageMaker reduces toil but does not remove design work. You still need disciplined data versioning, evaluation gates, cost controls, and clear ownership boundaries.
\item Use VPC attachment, KMS encryption, and separate execution roles when the platform handles regulated or proprietary data.
\item Choose SageMaker because the managed abstractions help, not because it is fashionable. Simple workloads may be cheaper elsewhere, while complex governed ML often benefits from it strongly.
\end{itemize}

\section*{What To Remember}
\begin{itemize}
\item SageMaker is a platform layer built on top of core AWS services, not a separate universe.
\item Training, deployment, and pipeline concerns should be understood independently even if SageMaker can connect them.
\item For exams and interviews, focus on when SageMaker reduces operational burden versus when EC2 or containers give needed control.
\end{itemize}

\end{document}
